% !TeX root = ../main.tex
% Fonti lette sulle immagini: Lezione 3.pdf, pp. 5--7;
% Lezione 4.pdf, pp. 2--14; Lezione 5.pdf, pp. 2--11.
% Precisazione rispetto alla formulazione sintetica del manoscritto:
% il gruppo moltiplicativo dei razionali esclude lo zero.
\chapter{Sistemi numerici e completezza}
\label{cap:sistemi-numerici}
\fonti{Lezioni 3, 4 e 5}

\section{Gli interi e la struttura di gruppo}
% Fonte: Lezione 3, p.~5

Gli interi estendono i naturali aggiungendo i numeri negativi:
\[
  \Z=\{0,\pm1,\pm2,\pm3,\ldots\}.
\]
Il vantaggio fondamentale è che ogni elemento ammette un opposto.
Nei naturali, per esempio, non esiste una soluzione di $1+y=0$;
negli interi la soluzione è $y=-1$.

La somma è un'operazione interna a $\Z$, cioè sommando due interi
si ottiene ancora un intero, e soddisfa:
\begin{enumerate}
  \item \emph{associatività}: $(x+y)+z=x+(y+z)$ per ogni $x,y,z\in\Z$;
  \item \emph{esistenza dell'elemento neutro}: $x+0=0+x=x$ per ogni $x\in\Z$;
  \item \emph{esistenza dell'opposto}: per ogni $x\in\Z$ esiste $-x\in\Z$
    tale che $x+(-x)=(-x)+x=0$.
\end{enumerate}
Queste proprietà esprimono che $(\Z,+)$ è un \emph{gruppo}.
Poiché vale anche la commutatività $x+y=y+x$, si tratta di un
gruppo \emph{commutativo}, detto anche \emph{abeliano}.

\section{I razionali come campo ordinato}
% Fonte: Lezione 3, p.~6

I numeri razionali sono quelli rappresentabili mediante una frazione
di interi con denominatore non nullo:
\[
  \Q=\left\{\frac pq:p,q\in\Z,\ q\ne0\right\}.
\]
Frazioni diverse possono rappresentare lo stesso numero:
$p/q=r/s$ se e solo se $ps=rq$, con $q,s\ne0$. Questo distingue
il numero razionale dalla particolare scrittura usata per esprimerlo.

La somma e il prodotto di razionali sono ancora razionali.
La somma rende $(\Q,+)$ un gruppo abeliano, come accade per gli interi.
Il prodotto rende invece $(\Q\setminus\{0\},\cdot)$ un gruppo abeliano:
l'elemento neutro è $1$ e l'inverso moltiplicativo di un razionale
non nullo $p/q$ è $q/p$. Lo zero è escluso dal gruppo moltiplicativo
perché non ha inverso. Inoltre il prodotto è distributivo rispetto
alla somma:
\[
  x(y+z)=xy+xz.
\]
Queste proprietà fanno di $(\Q,+,\cdot)$ un \emph{campo}.

L'ordine usuale $\le$ è totale anche su $\Q$ ed è compatibile con
le operazioni: sommare lo stesso numero conserva le disuguaglianze,
e il prodotto di due numeri non negativi è non negativo. Perciò
si parla di \emph{campo ordinato}.

\section{La proprietà di Archimede nei razionali}
% Fonte: Lezione 3, p.~6

Per ogni coppia di razionali $x,y>0$ esiste un naturale $n$ tale che
\[
  nx>y.
\]
Questa è la \emph{proprietà di Archimede}: per quanto piccolo sia
$x>0$ rispetto a $y$, un multiplo intero di $x$ riesce a superare $y$.

Nei razionali la proprietà si può verificare direttamente. Scriviamo
$y/x=p/q$, con $p,q$ interi positivi. Poiché $q\ge1$, si ha
$p/q\le p<p+1$. Basta allora scegliere il naturale $n=p+1$:
$n>y/x$ e, moltiplicando per $x>0$, si ottiene $nx>y$.

\section{La densità fra razionali}
% Fonte: Lezione 3, p.~7

Fra due razionali distinti esistono infiniti altri razionali:
\[
  \forall x,y\in\Q,\quad x<y
  \quad\Longrightarrow\quad
  \text{esistono infiniti }z\in\Q\text{ con }x<z<y.
\]
Per trovare un primo esempio basta la media $(x+y)/2$, che è ancora
razionale e si trova strettamente fra $x$ e $y$. Per verificare anche
che i valori possibili sono infiniti, consideriamo
\[
  z_m=x+\frac{y-x}{m+1},\qquad m\in\N^*.
\]
Sono tutti razionali, tutti distinti e tutti interni a $(x,y)$:
infatti $0<(y-x)/(m+1)<y-x$. A differenza dei naturali o degli interi,
i razionali non hanno quindi due elementi consecutivi.

\section{La numerabilità dei razionali}
% Fonte: Lezione 3, p.~7

Un insieme infinito si dice \emph{numerabile} quando i suoi elementi
si possono mettere in corrispondenza biunivoca con i naturali:
equivalentemente, si possono elencare senza omissioni e senza
ripetizioni in una lista indicizzata da $\N$.
Nonostante la densità appena vista, $\Q$ è numerabile.

Per costruire una lista cominciamo dai razionali positivi. Disponiamo
le frazioni $p/q$, con $p,q\in\N^*$, in una griglia: il numeratore
cambia lungo le colonne e il denominatore lungo le righe.
Percorriamo le diagonali di somma $p+q$ costante, partendo da $p+q=2$,
poi $3$, poi $4$, e così via. Su ciascuna diagonale prendiamo i
numeratori in ordine crescente.

\begin{schema}
\begin{tikzpicture}[x=1.45cm,y=.9cm]
  \foreach \numeratore in {1,...,4} {
    \node[black!65] at ({\numeratore-1},.7) {$p=\numeratore$};
    \foreach \denominatore in {1,...,4} {
      \node[fill=white,inner sep=4pt] (f-\numeratore-\denominatore)
        at ({\numeratore-1},{1-\denominatore})
        {$\dfrac{\numeratore}{\denominatore}$};
    }
  }
  \foreach \denominatore in {1,...,4} {
    \node[anchor=east,black!65] at (-.55,{1-\denominatore}) {$q=\denominatore$};
  }
  \foreach \somma in {3,4,5} {
    \pgfmathtruncatemacro{\massimo}{\somma-2}
    \foreach \numeratore in {1,...,\massimo} {
      \pgfmathtruncatemacro{\denominatore}{\somma-\numeratore}
      \pgfmathtruncatemacro{\successivo}{\numeratore+1}
      \pgfmathtruncatemacro{\precedente}{\denominatore-1}
      \draw[coloreschema,-{Stealth[length=1.7mm]},shorten <=2pt,shorten >=2pt]
        (f-\numeratore-\denominatore) -- (f-\successivo-\precedente);
    }
  }
  \draw[coloreschema,line width=.9pt] (.86,-1.21) -- (1.14,-.79);
\end{tikzpicture}
\par\smallskip
Le frecce percorrono le prime diagonali. La frazione $2/2$ si scarta,
perché rappresenta lo stesso numero di $1/1$, già incontrato.
\end{schema}

La lettura comincia quindi con
\[
  \frac11,\quad\frac12,\quad\frac21,\quad\frac13,\quad
  \frac22,\quad\frac31,\quad\frac14,\quad\ldots
\]
e si eliminano ogni volta i numeri già comparsi. Dopo questa
eliminazione, l'inizio della lista è
\[
  1,\quad\frac12,\quad2,\quad\frac13,\quad3,\quad
  \frac14,\quad\frac23,\quad\frac32,\quad4,\quad\ldots
\]
Ogni razionale positivo compare, perché una sua rappresentazione
$p/q$ si trova sulla diagonale $p+q$, raggiunta dopo un numero finito
di passi. Il procedimento non si esaurisce, dato che i razionali
positivi comprendono tutti i naturali positivi. Chiamando la lista
ottenuta $r_1,r_2,r_3,\ldots$, otteniamo infine una lista di tutti
i razionali, senza ripetizioni, nell'ordine
\[
  0,\ r_1,\ -r_1,\ r_2,\ -r_2,\ r_3,\ -r_3,\ldots
\]
e quindi una corrispondenza biunivoca con $\N$. Densità e numerabilità
descrivono dunque due aspetti differenti: la prima riguarda ciò che
si trova fra due numeri, la seconda la possibilità di elencare
tutti gli elementi dell'insieme.

% !TeX root = ../main.tex
% Lezione 4.pdf, pp. 2--5. Le precisazioni sulle sezioni e sui quantificatori
% completano la formulazione sintetica del manoscritto.
\section{I limiti dei razionali: l'irrazionalità della radice di due}
% Fonte: Lezione 4, p.~2

I razionali permettono di eseguire le quattro operazioni, con la sola
esclusione della divisione per zero, ma non risolvono ogni equazione
algebrica. L'esempio fondamentale è $x^2=2$: la sua soluzione positiva,
indicata con $\sqrt2$, non è razionale.

\begin{proposizione}
Non esiste $d\in\Q$ tale che $d^2=2$; in particolare $\sqrt2\notin\Q$.
\end{proposizione}
\begin{proof}
Supponiamo per assurdo che $d=p/q$, con $p,q\in\Z$, $q\ne0$ e
la frazione ridotta ai minimi termini, soddisfi $d^2=2$. Si avrebbe
\[
  p^2=2q^2.
\]
Quindi $p^2$ è pari e anche $p$ deve essere pari: infatti il quadrato
di un dispari $2h+1$ è $4h(h+1)+1$, che è dispari. Scriviamo allora
$p=2k$. Sostituendo nella relazione precedente otteniamo
\[
  4k^2=2q^2,\qquad q^2=2k^2.
\]
Anche $q^2$, e dunque $q$, è pari. Ma allora $p$ e $q$ sono entrambi
divisibili per $2$, in contraddizione con la scelta di una frazione
ridotta ai minimi termini.
\end{proof}

\section{I reali e l'assioma di Dedekind}
% Fonte: Lezione 4, pp.~3--4

Per descrivere la retta numerica si introduce il campo dei numeri reali
$\R$, che contiene $\Q$. La somma e il prodotto sono operazioni interne
e l'ordine $\le$ è totale. Le proprietà algebriche e di compatibilità
con l'ordine sono le seguenti:
\begin{enumerate}
  \item la somma è associativa e commutativa e ha elemento neutro $0$;
  \item ogni $x\in\R$ ha un opposto $-x$, con $x+(-x)=0$;
  \item il prodotto è associativo e commutativo e ha elemento neutro $1$,
    con $1\ne0$;
  \item ogni $x\ne0$ ha un reciproco $x^{-1}$, con $xx^{-1}=1$;
  \item il prodotto è distributivo rispetto alla somma:
    $x(y+z)=xy+xz$;
  \item se $a\le b$, allora $a+c\le b+c$ per ogni $c\in\R$;
  \item se $a\le b$ e $c\ge0$, allora $ac\le bc$.
\end{enumerate}
Queste proprietà fanno di $\R$ un campo ordinato. La proprietà che lo
distingue da $\Q$ è la \emph{completezza}, espressa qui dall'assioma
di Dedekind.

\begin{assioma}[Dedekind]
\label{ass:dedekind}
Siano $A,B\subseteq\R$ non vuoti e tali che $a\le b$ per ogni
$a\in A$ e ogni $b\in B$. Allora esiste $c\in\R$ per cui
\[
  a\le c\le b\qquad\text{per ogni }a\in A\text{ e ogni }b\in B.
\]
\end{assioma}

Il numero $c$ separa i due insiemi: si trova a destra di tutti gli
elementi di $A$ e a sinistra di tutti quelli di $B$. L'ordine dei
quantificatori richiede un valore di $c$ che funzioni
contemporaneamente per tutte le coppie:
\[
  \exists c\in\R\ \forall a\in A\ \forall b\in B:\ a\le c\le b.
\]
L'assioma afferma l'esistenza, senza richiedere che il separatore
sia determinato in modo univoco. Per esempio,
se $A=\{0\}$ e $B=\{1\}$, ogni $c\in[0,1]$ è un separatore.

\subsection{Perché l'assioma fallisce nei razionali}

Consideriamo gli insiemi
\[
  A=\{x\in\Q:x\ge0,\ x^2\le2\},\qquad
  B=\{x\in\Q:x\ge0,\ x^2\ge2\}.
\]
Sono entrambi non vuoti: $1\in A$ e $2\in B$. Inoltre $a\le b$
per ogni $a\in A$ e $b\in B$, perché sono non negativi e
$a^2\le2\le b^2$. Tuttavia nessun razionale può separare $A$ e $B$.

Per verificarlo, un eventuale separatore razionale $c$ dovrebbe
soddisfare $1\le c\le2$. Se $c^2<2$, scegliamo un razionale
\[
  0<h<\min\left\{1,\frac{2-c^2}{2c+1}\right\}.
\]
Allora $(c+h)^2<c^2+(2c+1)h<2$, quindi $c+h\in A$ ma $c+h>c$:
il separatore non sarebbe a destra di tutti gli elementi di $A$.
Se invece $c^2>2$, scegliamo un razionale
\[
  0<h<\min\left\{1,\frac{c^2-2}{2c}\right\}.
\]
Si ha $c-h>0$ e $(c-h)^2>c^2-2ch>2$, perciò $c-h\in B$ ma
$c-h<c$, ancora una contraddizione. Rimane soltanto la possibilità
$c^2=2$, esclusa dall'irrazionalità appena dimostrata. In $\R$ il
separatore è il punto $\sqrt2$ che manca a $\Q$.

\subsection{L'idea delle sezioni di Dedekind}
% Fonte: Lezione 4, p.~5

È possibile costruire i reali a partire dai razionali mediante le
\emph{sezioni di Dedekind}. Si considera una coppia $(A,B)$ di sottoinsiemi non vuoti
di $\Q$ che soddisfi
\[
  A\cap B=\varnothing,\qquad A\cup B=\Q,\qquad
  a\le b\quad\forall a\in A,\ \forall b\in B.
\]
La coppia descrive un \emph{taglio} della retta razionale in una parte
inferiore e una superiore. Poiché le due parti sono disgiunte,
in realtà $a<b$ per ogni $a\in A$, $b\in B$.

Per rendere la rappresentazione univoca occorre fissare una convenzione
sull'estremo: scegliamo che $A$ non abbia massimo. Infatti, per un
razionale $r$, sia la coppia
$\bigl(\{q:q<r\},\{q:q\ge r\}\bigr)$ sia la coppia
$\bigl(\{q:q\le r\},\{q:q>r\}\bigr)$ soddisfano le tre condizioni
precedenti, pur descrivendo lo stesso punto. La convenzione sceglie
la prima. Con questa precisazione, tali sezioni possono essere usate,
con opportune operazioni e ordine, per costruire un modello di $\R$.
Un taglio può rappresentare anche un punto irrazionale: in quel caso
nessuna delle due parti contiene un estremo razionale nel punto
di separazione.

% !TeX root = ../main.tex
% Lezione 4.pdf, pp. 6--14. Le prove dei criteri operativi e della
% proprietà duale dell'infimo rendono espliciti i passaggi della fonte.
\section{Intervalli e retta reale estesa}
% Fonte: Lezione 4, pp.~6--8

\begin{definizione}[Intervallo]
Un sottoinsieme non vuoto $I\subseteq\R$ è un \emph{intervallo} se,
ogni volta che $x,y\in I$ e $x<z<y$, anche $z$ appartiene a $I$.
\end{definizione}

In un intervallo si trovano quindi tutti i punti compresi fra due
suoi elementi. Per estremi finiti $a<b$ le principali possibilità sono
\begin{align*}
  [a,b]&=\{x\in\R:a\le x\le b\},&
  (a,b)&=\{x\in\R:a<x<b\},\\
  [a,b)&=\{x\in\R:a\le x<b\},&
  (a,b]&=\{x\in\R:a<x\le b\}.
\end{align*}
I primi due sono detti rispettivamente intervallo \emph{chiuso} e
\emph{aperto}; negli altri due è incluso un solo estremo. Anche
$[a,a]=\{a\}$ è un intervallo.\footnote{La definizione richiede
che un intervallo sia non vuoto. Altre convenzioni ammettono anche
l'insieme vuoto come intervallo.}

Per $I=[a,b]$, i numeri $a,b$ sono gli \emph{estremi}, $b-a$ è
l'\emph{ampiezza}, $(b-a)/2$ il \emph{raggio} e $(a+b)/2$ il
\emph{centro}. Le stesse espressioni individuano ampiezza e centro
anche se uno o entrambi gli estremi non sono inclusi.

Per descrivere intervalli illimitati introduciamo i simboli
$-\infty$ e $+\infty$, ordinati in modo che
\[
  -\infty<x<+\infty\qquad\forall x\in\R.
\]
La \emph{retta reale estesa} è
$\overline\R=\R\cup\{-\infty,+\infty\}$. Le semirette della retta
reale si scrivono, per esempio,
\begin{align*}
  (-\infty,c]&=\{x\in\R:x\le c\},&
  (-\infty,c)&=\{x\in\R:x<c\},\\
  [a,+\infty)&=\{x\in\R:x\ge a\},&
  (a,+\infty)&=\{x\in\R:x>a\}.
\end{align*}
In queste scritture $\pm\infty$ non appartengono agli intervalli,
che sono sottoinsiemi di $\R$; anche tutta $\R=(-\infty,+\infty)$
è un intervallo.

\section{Limitatezza, maggioranti e minoranti}
% Fonte: Lezione 4, pp.~8--10

\begin{definizione}
Un insieme non vuoto $A\subseteq\R$ è \emph{superiormente limitato}
se esiste $M\in\R$ tale che $a\le M$ per ogni $a\in A$.
Un tale numero è un \emph{maggiorante} di $A$.
È \emph{inferiormente limitato} se esiste $m\in\R$ tale che
$m\le a$ per ogni $a\in A$; un tale numero è un \emph{minorante}.
\end{definizione}

L'insieme è \emph{limitato} se lo è sia superiormente sia inferiormente:
\[
  \exists m,M\in\R:\quad m\le a\le M\quad\forall a\in A.
\]
Equivalentemente, esiste $N\ge0$ per cui $|a|\le N$ per ogni
$a\in A$. Infatti, da due estremi $m,M$ basta scegliere
$N=\max\{|m|,|M|\}$; viceversa, $|a|\le N$ equivale a
$-N\le a\le N$. Un insieme che non è limitato si dice
\emph{illimitato}; si può anche specificare la direzione in cui
manca un limite, parlando di insieme superiormente o inferiormente
illimitato.

Maggioranti e minoranti non devono appartenere all'insieme e non
sono determinati in modo univoco. Per esempio, sono possibili scelte
anche molto lontane dagli estremi:
\begin{itemize}
  \item per $A=[0,1)$ si possono prendere $m=-38$ e $M=251$;
  \item per $B=\{1/n:n\in\N^*\}$ si possono prendere
    $m=-1719{,}3$ e $M=16$, perché $0<1/n\le1$;
  \item per $C=\N$, $0$ è un minorante ma non esistono maggioranti,
    come seguirà dal principio di Archimede
    (teorema~\ref{teo:archimede-reali}).
\end{itemize}

Per $A=[0,3)$, in particolare, sono maggioranti $3$, $5$, $11$,
$197$ e, più in generale, tutti e soli i reali $M\ge3$.
Sono minoranti $0$, $-13$, $-\pi$, $-171$ e, in generale, tutti e soli
i reali $m\le0$. La condizione vale per ogni elemento di $A$,
anche se $3\notin A$.

\section{Massimo e minimo}
% Fonte: Lezione 4, p.~11

\begin{definizione}
Un \emph{massimo} di un insieme non vuoto $A\subseteq\R$ è un
maggiorante che appartiene ad $A$; un \emph{minimo} è un minorante
che appartiene ad $A$. Si scrive
\begin{align*}
  M=\max A&\quad\Longleftrightarrow\quad
    M\in A\text{ e }a\le M\ \forall a\in A,\\
  m=\min A&\quad\Longleftrightarrow\quad
    m\in A\text{ e }m\le a\ \forall a\in A.
\end{align*}
\end{definizione}

Se esistono, massimo e minimo sono unici. Per esempio, se $M_1$ e
$M_2$ fossero entrambi massimi, l'appartenenza di ciascuno ad $A$
imporrebbe $M_1\le M_2$ e $M_2\le M_1$, quindi $M_1=M_2$.
Per il minimo si ragiona allo stesso modo.

L'intervallo $[0,3)$ ha minimo $0$ e non ha massimo: per ogni
$a\in[0,3)$ il punto $(a+3)/2$ è ancora nell'intervallo ed è
strettamente maggiore di $a$. Quindi l'essere limitato non assicura
che un insieme possieda un massimo o un minimo.

\section{Estremo superiore ed estremo inferiore}
% Fonte: Lezione 4, p.~12

\begin{definizione}[Supremo e infimo]
Il numero $S\in\R$ è l'\emph{estremo superiore}, o \emph{supremo},
di $A$ se è il più piccolo dei suoi maggioranti. Il numero $s\in\R$
è l'\emph{estremo inferiore}, o \emph{infimo}, se è il più grande
dei suoi minoranti. In formule,
\begin{align*}
  S=\sup A&=\min\{M\in\R:a\le M\ \forall a\in A\},\\
  s=\inf A&=\max\{m\in\R:m\le a\ \forall a\in A\},
\end{align*}
quando tali estremi esistono.
\end{definizione}

Anche questi valori, quando esistono, sono unici. A differenza del
massimo e del minimo, possono non appartenere ad $A$. Se un massimo
esiste, è anche il più piccolo maggiorante: essendo un elemento di
$A$, è minore o uguale a qualunque altro maggiorante. Pertanto
\[
  \max A=\sup A\quad\text{se il massimo esiste},\qquad
  \min A=\inf A\quad\text{se il minimo esiste}.
\]
Viceversa, se $\sup A\in A$ esso è il massimo, e se $\inf A\in A$
esso è il minimo.

\subsection{Il teorema di esistenza degli estremi}
% Fonte: Lezione 4, pp.~13--14

\begin{teorema}[Esistenza del supremo e dell'infimo]
\label{teo:esistenza-estremi}
Ogni sottoinsieme non vuoto di $\R$ superiormente limitato ha
un estremo superiore reale. Ogni sottoinsieme non vuoto di $\R$
inferiormente limitato ha un estremo inferiore reale.
\end{teorema}
\begin{proof}
Sia $A\subseteq\R$ non vuoto e superiormente limitato. L'insieme
dei suoi maggioranti
\[
  \mathcal M_A=\{M\in\R:a\le M\ \forall a\in A\}
\]
è non vuoto. Per definizione, ogni elemento di $A$ è minore o
uguale a ogni elemento di $\mathcal M_A$. Applichiamo l'assioma
di Dedekind ai due insiemi $A$ e $\mathcal M_A$: esiste $S\in\R$
tale che
\[
  a\le S\le M\qquad\forall a\in A,\ \forall M\in\mathcal M_A.
\]
La prima disuguaglianza dice che $S$ è un maggiorante di $A$;
la seconda che è minore o uguale a ogni maggiorante. Dunque
$S=\min\mathcal M_A=\sup A$.

Per l'infimo si può applicare lo stesso ragionamento all'insieme
$-A=\{-a:a\in A\}$. Se $A$ è inferiormente limitato, $-A$ è
superiormente limitato e ha supremo $T$. Cambiando segno alle
disuguaglianze si ottiene che $-T$ è il più grande minorante di
$A$, ossia $\inf A=-\sup(-A)$.
\end{proof}

Il teorema è una conseguenza della completezza. Sostituendo $\R$
con $\Q$ l'enunciato diventa falso: l'insieme
\[
  A=\{x\in\Q:x\ge0,\ x^2\le2\}
\]
è non vuoto e superiormente limitato, ma non ha supremo in $\Q$.
Un eventuale minimo maggiorante razionale sarebbe infatti un
separatore razionale fra gli insiemi $A$ e $B$ già considerati
nell'esempio sull'assioma di Dedekind. Considerando lo stesso
insieme come sottoinsieme di $\R$, il suo supremo è $\sqrt2$;
il suo infimo è $0$, che è anche minimo.

\subsection{Estremi infiniti}

Per gli insiemi non vuoti si estende la notazione ponendo
$\sup A=+\infty$ quando $A$ è superiormente illimitato e
$\inf A=-\infty$ quando è inferiormente illimitato.
Si tratta di valori nella retta estesa, non di massimi o minimi
reali. Per esempio, $\sup\N=+\infty$ e $\inf\N=\min\N=0$.

\section{La caratterizzazione operativa di supremo e infimo}
% Fonte: Lezione 4, p.~14; lezione 5, p.~2

\begin{proposizione}
\label{prop:criterio-estremi}
Sia $A\subseteq\R$ non vuoto. Un numero $S\in\R$ è il suo supremo
se e solo se valgono entrambe le condizioni
\[
  \begin{cases}
    a\le S&\forall a\in A,\\
    \forall\varepsilon>0\ \exists a\in A:\ a>S-\varepsilon.
  \end{cases}
\]
Analogamente, $s\in\R$ è l'infimo se e solo se
\[
  \begin{cases}
    s\le a&\forall a\in A,\\
    \forall\varepsilon>0\ \exists a\in A:\ a<s+\varepsilon.
  \end{cases}
\]
\end{proposizione}
\begin{proof}
Se $S=\sup A$, esso è un maggiorante. Se per qualche
$\varepsilon>0$ non esistesse $a\in A$ con $a>S-\varepsilon$,
tutti gli elementi sarebbero minori o uguali a $S-\varepsilon$:
quest'ultimo sarebbe un maggiorante più piccolo di $S$, assurdo.

Viceversa, supponiamo che le due condizioni valgano. $S$ è un
maggiorante e ogni $M<S$ non può esserlo: scegliendo
$\varepsilon=S-M>0$ si trova $a\in A$ con $a>M$. Quindi $S$ è
il più piccolo dei maggioranti.

Per l'infimo, se $s=\inf A$ e mancassero elementi sotto
$s+\varepsilon$, quest'ultimo sarebbe un minorante più grande.
Viceversa, se le due condizioni per $s$ valgono, ogni $m>s$ viene
escluso dai minoranti scegliendo $\varepsilon=m-s$ e trovando
$a\in A$ con $a<m$. Dunque $s$ è il più grande minorante.
\end{proof}

Il criterio traduce il ruolo di estremo in due verifiche separate:
essere un limite per tutti gli elementi e poter essere avvicinato
arbitrariamente dall'insieme. L'elemento $a$ può dipendere da
$\varepsilon$. Se l'estremo appartiene ad $A$, lo si può scegliere
per ogni $\varepsilon$; il criterio è particolarmente utile quando
l'estremo non è raggiunto.

\begin{schema}
\begin{tikzpicture}[x=1cm,y=1cm]
  \fill[coloreschema!12] (4.3,.52) rectangle (6.8,.78);
  \fill[coloreschema!12] (1.2,-.98) rectangle (3.7,-.72);
  \foreach \quota in {.65,-.85} {
    \draw[asse] (0,\quota) -- (8,\quota);
  }
  \node[anchor=east] at (-.2,.65) {$S=\sup A$};
  \node[anchor=east] at (-.2,-.85) {$s=\inf A$};
  \foreach \pos/\testo in {4.3/{S-\varepsilon},6.8/S} {
    \draw (\pos,.55) -- (\pos,.75);
    \node[below=6pt] at (\pos,.65) {$\testo$};
  }
  \foreach \pos/\testo in {1.2/s,3.7/{s+\varepsilon}} {
    \draw (\pos,-.95) -- (\pos,-.75);
    \node[below=6pt] at (\pos,-.85) {$\testo$};
  }
  \fill[coloreschema] (6.1,.65) circle[radius=2pt];
  \node[above=6pt,coloreschema] at (6.1,.65) {$a\in A$};
  \fill[coloreschema] (2.1,-.85) circle[radius=2pt];
  \node[above=6pt,coloreschema] at (2.1,-.85) {$a\in A$};
\end{tikzpicture}
\par\smallskip
Per ogni $\varepsilon>0$, la fascia $(S-\varepsilon,S]$ deve
contenere almeno un punto di $A$; analogamente per $[s,s+\varepsilon)$.
Il colore evidenzia la fascia di ricerca, senza affermare che tutti
i suoi punti appartengano ad $A$.
\end{schema}

% !TeX root = ../main.tex
% Lezione 5.pdf, pp. 2--11. La proprietà di Archimede precede gli esempi
% che la usano. La parte sulle funzioni, pp. 12--16, è nel capitolo 05.
% Corretto il passaggio del manoscritto su 1/n < -1+epsilon:
% per 0<epsilon<1 il secondo membro è negativo e non vi sono soluzioni.
\section{Il principio di Archimede nei reali}
% Fonte: Lezione 5, pp.~8--9

\begin{teorema}[Principio di Archimede]
\label{teo:archimede-reali}
Fissati $a,b\in\R$, con $a>0$, esiste $n\in\N$ tale che $na>b$.
\end{teorema}
\begin{proof}
\textbf{Riduzione al supremo dei multipli.} Fissiamo $a>0$ e consideriamo l'insieme
\[
  A=\{na:n\in\N\}.
\]
È non vuoto, perché $0\in A$. Dimostriamo che è superiormente illimitato. Se per assurdo fosse superiormente limitato, il teorema di esistenza degli estremi garantirebbe un numero reale $S=\sup A$.

\textbf{Il successivo supera il supremo.} La caratterizzazione operativa, con la scelta $\varepsilon=a>0$, darebbe un elemento $\bar n a\in A$ per cui
\[
  \bar n a>S-a.
\]
Sommando $a$ ai due membri otteniamo $(\bar n+1)a>S$. Ma $\bar n+1\in\N$, quindi anche $(\bar n+1)a\in A$, in contraddizione con il fatto che $S$ sia un maggiorante. L'insieme dei multipli è dunque superiormente illimitato: per ogni $b\in\R$ contiene un elemento $na>b$, come richiesto.
\end{proof}

\textbf{Conseguenze.}
\begin{itemize}
  \item Con $a=1$ si ottiene che i naturali superano qualunque reale: $\N$ è superiormente illimitato.
  \item I reciproci dei naturali positivi possono essere arbitrariamente piccoli:
  \[
    \forall\varepsilon>0\ \exists n\in\N^*:\quad \frac1n<\varepsilon.
  \]
  Basta infatti scegliere $n>1/\varepsilon$.
\end{itemize}
La completezza ha quindi una conseguenza concreta: non esiste un reale positivo minore o uguale a tutti i reciproci dei naturali positivi.

\addtocontents{toc}{\protect\Needspace{3\protect\baselineskip}}
\section{Esempi ed esercizi su supremo e infimo}
% Fonte: Lezione 5, pp.~2--7

\subsection{Un intervallo e un insieme con un punto isolato}
% Fonte: Lezione 5, p.~2

\begin{esempio}[Intervallo con estremo escluso]
Per $A=[0,1)$ si ha
\[
  \inf A=\min A=0,\qquad \sup A=1.
\]
\textbf{Minimo.} Il minimo è $0$, perché $0\in A$ e tutti gli elementi sono non negativi.

\textbf{Supremo non raggiunto.} Il numero $1$ è un maggiorante e, dato $\varepsilon>0$, il punto $a=1-\min\{\varepsilon/2,1/2\}$ appartiene ad $A$ e soddisfa $a>1-\varepsilon$. Il criterio operativo dimostra dunque che $1=\sup A$. Tuttavia $1\notin A$, quindi il massimo non esiste.
\end{esempio}

\begin{esempio}[Un massimo isolato]
Consideriamo
\[
  B=[-1,2]\cup\{5\}.
\]
In questo caso entrambi gli estremi sono raggiunti:
\[
  \inf B=\min B=-1,\qquad \sup B=\max B=5.
\]
Per verificare il criterio del supremo basta scegliere $a=5$ per ogni $\varepsilon>0$, poiché $5>5-\varepsilon$. Il punto che realizza il massimo può dunque essere isolato; la presenza dell'intero intervallo fra $2$ e $5$ non è richiesta.
\end{esempio}

\begin{schema}
\begin{tikzpicture}[x=1.45cm,y=1cm]
  \foreach \quota/\nome in {.7/A,-.8/B} {
    \draw[asse] (-1.5,\quota) -- (5.5,\quota);
    \node[anchor=east] at (-1.8,\quota) {$\nome$};
  }
  \draw[intervallo] (0,.7) -- (1,.7);
  \fill[coloreschema] (0,.7) circle[radius=2.5pt];
  \draw[estremo] (1,.7) circle[radius=2.5pt];
  \node[below=6pt] at (0,.7) {$0$};
  \node[below=6pt] at (1,.7) {$1$};
  \draw[intervallo] (-1,-.8) -- (2,-.8);
  \foreach \pos in {-1,2,5} {
    \fill[coloreschema] (\pos,-.8) circle[radius=2.5pt];
    \node[below=6pt] at (\pos,-.8) {$\pos$};
  }
\end{tikzpicture}
\par\smallskip In $A$ l'estremo superiore $1$ è escluso; in $B$ il punto $5$ è presente ed è il massimo, anche se è separato dall'intervallo.
\end{schema}

\subsection{I reciproci dei naturali positivi}
% Fonte: Lezione 5, pp.~3--4

\begin{esempio}[Reciproci dei naturali positivi]
Sia
\[
  A=\left\{\frac1n:n\in\N^*\right\}.
\]
La rappresentazione dei suoi punti suggerisce $\sup A=1$ e $\inf A=0$. Verifichiamo entrambe le affermazioni.

\textbf{Massimo e supremo.} Il numero $1$ appartiene ad $A$ per $n=1$ e maggiora tutti gli elementi, poiché $1/n\le1$ per ogni $n\ge1$. Dunque
\[
  1=\max A=\sup A.
\]
\textbf{Infimo non raggiunto.} Il numero $0$ non appartiene ad $A$: nessun reciproco di un naturale positivo è zero. Per dimostrare che è l'infimo usiamo le due condizioni operative. È un minorante, perché $1/n>0$; inoltre, per ogni $\varepsilon>0$, Archimede fornisce un naturale $n>1/\varepsilon$, e quindi
\[
  \frac1n<\varepsilon=0+\varepsilon.
\]
Ne segue $\inf A=0$. Poiché $0\notin A$, il minimo non esiste.
\end{esempio}

\subsection{Come scartare un candidato sbagliato per l'infimo}
% Fonte: Lezione 5, pp.~5--6

Per lo stesso insieme $A=\{1/n:n\in\N^*\}$ proviamo a proporre un candidato $k$ diverso da $0$.

\textbf{Candidato positivo: non è un minorante.} Se $k>0$, $k$ non è un minorante. Infatti esiste $n>1/k$, quindi $1/n<k$.

\textbf{Candidato negativo: manca l’avvicinabilità.} Se $k<0$, esso è un minorante, perché tutti gli elementi di $A$ sono positivi, ma non è il più grande: $0$ è un minorante maggiore di $k$. Lo stesso difetto si vede con il criterio operativo. Per il candidato $k=-1$, prendendo $0<\varepsilon<1$ avremmo bisogno di un elemento per cui
\[
  \frac1n<-1+\varepsilon<0,
\]
che è impossibile. In generale, per un qualunque $k<0$, basta scegliere $\varepsilon=-k/2$: allora $k+\varepsilon=k/2<0$ e nessun elemento di $A$ può trovarsi sotto questa soglia.

\subsection{Due esercizi: cubi traslati e unione di insiemi}
% Fonte: Lezione 5, p.~7

\begin{esercizio}[Cubi traslati]
Determinare supremo e infimo dell'insieme
\[
  A_1=\{n^3-1:n\in\N\}.
\]
Stabilire se esistono massimo e minimo.
\end{esercizio}
\begin{svolgimento}
Poiché $n^3\ge0$ per ogni $n\in\N$ e il valore $-1$ si ottiene per $n=0$, si ha $\inf A_1=\min A_1=-1$.

\textbf{Illimitatezza superiore.} L'insieme è superiormente illimitato: dato un qualunque $M\in\R$, scegliamo un naturale $n>\max\{1,M+1\}$. Allora $n^3-1\ge n-1>M$. Dunque, nella retta estesa,
\[
  \sup A_1=+\infty.
\]
L'illimitatezza superiore esclude l'esistenza di un massimo.
\end{svolgimento}

\begin{esercizio}[Unione di due insiemi]
Determinare supremo, infimo ed eventuali massimo e minimo dell'insieme $A_2=\{x\in\R:x^2<1\}\cup\{x\in\R:(x-3)^2=0\}$.
\end{esercizio}
\begin{svolgimento}
\textbf{Descrizione dell'insieme.} Si ha
\[
  A_2=\{x\in\R:x^2<1\}
      \cup\{x\in\R:(x-3)^2=0\}
     =(-1,1)\cup\{3\}.
\]
Il numero $3$ appartiene ad $A_2$ e maggiora tutti gli elementi, quindi $\sup A_2=\max A_2=3$.

\textbf{Infimo non raggiunto.} Il numero $-1$ è un minorante. Dato $\varepsilon>0$, scegliamo
\[
  a=-1+\frac{\varepsilon}{1+\varepsilon}.
\]
Il punto $a$ appartiene a $(-1,0)\subseteq A_2$ e soddisfa $a<-1+\varepsilon$, perché $0<\varepsilon/(1+\varepsilon)<\varepsilon$. Perciò $\inf A_2=-1$; il minimo non esiste, perché $-1\notin A_2$.
\end{svolgimento}

% Mantiene insieme la dimostrazione e il disegno della densità.
\Needspace{36\baselineskip}
\section{La densità dei razionali nella retta reale}
% Fonte: Lezione 5, pp.~10--11

Fra due razionali distinti si trovano altri razionali. Il risultato seguente estende la proprietà a qualunque intervallo aperto con estremi reali.

\begin{teorema}[Densità di $\Q$ in $\R$]
\label{teo:densita-razionali}
Se $a,b\in\R$ e $a<b$, esiste $r\in\Q$ tale che $a<r<b$.
\end{teorema}
\begin{proof}
\textbf{Dilatazione dell’intervallo.} Poiché $b-a>0$, il principio di Archimede permette di scegliere $q\in\N^*$ tale che
\[
  q(b-a)>1.
\]
L'intervallo $(qa,qb)$ ha dunque ampiezza maggiore di $1$ e contiene un intero $p$.

\textbf{Scelta dell’intero.} Per ogni reale $t$ esiste un intero $p$ con $t<p\le t+1$. Infatti, scegliamo un naturale $N>-t$ e consideriamo il minimo naturale $m>t+N$, che esiste per Archimede e per il principio del minimo intero. Si ha $m\ge1$ e $m-1\le t+N$, altrimenti $m$ non sarebbe il minimo. Ponendo $p=m-N$ si ottiene proprio $t<p\le t+1$.

\textbf{Ritorno alla scala iniziale.} Applicando il fatto appena verificato a $t=qa$, troviamo
\[
  qa<p\le qa+1<qb.
\]
Dividendo per $q>0$ segue $a<p/q<b$. Il numero $r=p/q$ è razionale, perché $p\in\Z$ e $q\in\N^*$, e appartiene all'intervallo richiesto.
\end{proof}

\begin{schema}
\begin{tikzpicture}[x=1.6cm,y=1cm]
  \draw[asse] (-.3,0) -- (4.4,0);
  \draw[intervallo] (1.35,0) -- (3.25,0);
  \draw[estremo] (1.35,0) circle[radius=2.5pt];
  \draw[estremo] (3.25,0) circle[radius=2.5pt];
  \node[above=8pt,coloreschema] at (1.35,0) {$qa$};
  \node[above=8pt,coloreschema] at (3.25,0) {$qb$};
  \foreach \pos/\testo in {0/{p-2},1/{p-1},2/p,3/{p+1},4/{p+2}} {
    \draw[black!55] (\pos,-.08) -- (\pos,.08);
    \node[below=7pt] at (\pos,0) {$\testo$};
  }
  \fill[coloreschema] (2,0) circle[radius=2.5pt];
\end{tikzpicture}
\par\smallskip Dilatando l'intervallo con un fattore intero $q$ si ottiene un'ampiezza maggiore di $1$, che assicura la presenza di un intero $p$. Tornando alla scala iniziale si trova il razionale $p/q$.
\end{schema}

\textbf{Infinità dei punti razionali.} I razionali in $(a,b)$ sono infiniti: trovato un primo razionale $r_1$ nell'intervallo, si può cercarne un secondo in $(a,r_1)$, un terzo fra $a$ e il secondo, e continuare. Ogni intervallo aperto non vuoto della retta incontra quindi $\Q$, anche quando i suoi estremi sono irrazionali.

